% =============================================================================
% GREEDY BORUTA - ECML-PKDD SUBMISSION WIREFRAME
% =============================================================================
% Format: Springer LNCS (Lecture Notes in Computer Science)
% Max length: 16 pages including references (Research Track)
% Review: Double-blind
% =============================================================================

%Nicolas, algumas observacoes sobre o paper:

% 1 - Introduzimos o conceito de weakly and strongly confirmed features. Isso tem impacto no output da biblitoca greedy boruta, mas nao tem muito impacto no paper em si.

% 2 - Introduzimos a analise de FPR pra quando o numero de shadows eh r <= p. Eu nao mexi nas secoes de resultados do paper, porque todas as secoes usam p shadow features. No entanto, eu criei uma nova secao (e esta vazia) no apendice B onde podemos colocar resultados do GB usando r shadow features versus o vanilla boruta e o GB com p shadow features. Estava pensando em pegar os datasets os 4 datasets com os maiores numeros de features e meio que escrever uma secao bem resumida dos resultados no estilo da secao q esta no paper. O q acha? Se achar interessante, compartilha os arquivos que geram os resultados da secao de dados reais que eu me viro.

% 3 - Estava pensando em mover a secao "Reducing the Number of Shadow Features" pro Apendice B. A ideia eh legal, mas nao sei se faz tanto sentido deixar essa secao no meio do paper. Me diz o que vc acha.

% 4 - Eu fiz varias alteracoes, algumas maiores e outras menores. Dessa vez eu nao aceitei, entao o paper ta cheio que alteracoes pra vc ver onde eu mexi exatamente.

% 5 - Remover o titulo “Undecidedness & …” da figura 5. A legenda ja diz a mesma coisa.

% Idem pra Figura 4. A legenda e o titulo dizem a mesma coisa.

% Remover o titutlo “Ground-Truth metrics by alpha” da figura 2.

% 6 - Pra mandarmos pro Chris, so preciso que vc de uma olhada no paper e resolva os comentarios. Todos os pontos acima podem ser resolvido apos o envio do paper ao Chris, pq ele nao vai conseguir ler o paper de imediato. 

\documentclass[runningheads]{llncs}

% =============================================================================
% PACKAGES
% =============================================================================
\usepackage{graphicx}
\let\proof\relax
\let\endproof\relax
\usepackage{amsthm}
\usepackage{amsmath,amssymb}
\usepackage{esvect}
\usepackage{algorithm}
\usepackage{algorithmic}
\usepackage{booktabs}
\usepackage{multirow}
\usepackage{xcolor}
\usepackage{hyphenat}
\usepackage{subcaption}
\usepackage{xurl}
\usepackage{hyperref}
\usepackage{cleveref}
\usepackage{microtype}

% \usepackage[htt]{hyphenat}
% DRAFT ONLY: forces figures to appear exactly where declared in source (no floating).
% TO REVERT FOR SUBMISSION: remove the float package import below and change all
% \begin{figure}[H] back to \begin{figure}[t] (or [htbp] for more placement flexibility).
\usepackage{float}
% DRAFT ONLY: cap figure height so images don't fill an entire page and leave blank gaps.
% TO REVERT FOR SUBMISSION: remove the two \setlength lines below, and remove the
% height=... keepaspectratio arguments from each \includegraphics call in section 5.
\setlength{\abovecaptionskip}{4pt}
\setlength{\belowcaptionskip}{0pt}

% For wireframe TODO notes - REMOVE BEFORE SUBMISSION
\usepackage[colorinlistoftodos,textsize=small]{todonotes}
\newcommand{\writeme}[1]{\todo[inline,color=yellow!40]{WRITE: #1}}
\newcommand{\addcite}[1]{\todo[inline,color=blue!20]{CITE: #1}}
\newcommand{\addfigure}[1]{\todo[inline,color=green!20]{FIGURE: #1}}
\newcommand{\addtable}[1]{\todo[inline,color=orange!20]{TABLE: #1}}
\newcommand{\decision}[1]{\todo[inline,color=red!20]{DECISION: #1}}

\newtheorem{assumption}{Assumption}

% =============================================================================
% TITLE AND AUTHORS
% =============================================================================
\begin{document}

\title{Trading statistical conservatism for speed: A `greedy' modification to Boruta}
% Alternative titles:
% - Trading Statistical Conservatism for Speed: A Greedy Modification to Boruta
% - Fast All-Relevant Feature Selection with Provable Guarantees

\titlerunning{Greedy Boruta}

% For double-blind submission, use anonymous authors
\author{Nicolas Vana, Estevão Prado}
% \author{}
\authorrunning{Anonymous}
% \institute{Anonymous Institution}

% For camera-ready version, uncomment and fill:
\author{Nicolas Vana \and Estevão Prado}
\authorrunning{Vana and Prado}
\institute{Bradesco bank, São Paulo, Brazil}
% \email{nicolas.v.santos@bradesco.com.br}
% \email{nicolas.vana@bradesco.com.br}
% \and Institution 2, City, Country}

\maketitle

% =============================================================================
% ABSTRACT (~150-200 words)
% =============================================================================

\begin{abstract}
The Boruta algorithm is the most widely adopted method for all-relevant feature selection, but its computational cost and lack of convergence guarantees limit its applicability in modern high-dimensional settings, where the number of features $p$ frequently exceeds $10^4$. We propose Greedy Boruta, a modification that replaces Vanilla Boruta's binomial confirmation test with a single-hit rule: a feature is confirmed if it ever exceeds the maximum shadow importance in any iteration. Despite its simplicity, this change yields four formal guarantees: convergence in exactly $K = \lceil \log_2(p/\alpha) \rceil$ iterations, feature recall at least as high as Vanilla Boruta, triviality of FDR correction under the greedy criterion and a per-feature false positive rate bounded above by $K/(p+1)$, which decreases as dimensionality grows. The deterministic iteration count eliminates all termination-related hyperparameters and enables embarrassingly parallel execution, implemented in \texttt{greedyboruta} (available on PyPI). We further introduce a weak/strong confirmation framework that distinguishes features satisfying the original Boruta criterion from those selected exclusively by the greedy criterion, and show that the number of shadow features can be reduced from $p$ to $r \leq p$, providing an explicit trade-off between computational cost and false positive control. On a synthetic benchmark with known ground truth and six real-world high-dimensional datasets, Greedy Boruta achieves 1.8 to 66.2 times speedups while not performing significantly worse than Vanilla Boruta in downstream model performance across a set of real-world datasets used as benchmarks.
\keywords{Feature selection \and Boruta algorithm \and All-relevant features \and Computational efficiency \and Parallelization}
\end{abstract}

\section{Introduction}
\label{sec:introduction}

Feature selection is a critical preprocessing step in high-dimensional machine learning pipelines, reducing computational cost, improving model interpretability and mitigating overfitting. Methods for feature selection broadly fall into two paradigms. \emph{Minimal-optimal} methods, such as Lasso~\cite{tibshirani1996lasso}, Ridge~\cite{hoerl1970ridge}, Recursive Feature Elimination~\cite{guyon2002rfe} and mRMR~\cite{peng2005mrmr}, seek the smallest subset of features that maximizes predictive performance. \emph{All-relevant} methods instead aim to identify every feature that carries information about the target variable, regardless of redundancy. The all-relevant paradigm is particularly valuable in scientific discovery~\cite{rak2019evaluation}, biomarker identification~\cite{mardian2023ovarian} and causal inference~\cite{pellet2008markov}, where discarding a relevant feature represents a permanent loss of information that no downstream analysis can recover.

The Boruta algorithm~\cite{kursa2010feature,kursa2010boruta} is the most widely adopted method for all-relevant feature selection. Its statistical framework, based on comparing each feature's importance against permuted ``shadow'' copies that represent the maximum achievable importance by random noise, provides a principled mechanism for distinguishing signal from noise. Since its introduction, Boruta has been adopted across diverse domains including bioinformatics~\cite{rak2019evaluation,mardian2023ovarian}, finance~\cite{pinto2024forecasting,li2024stock} and natural language processing~\cite{sallah2024fake,tan2024mosquito}, and several variants have been proposed. For instance, BorutaPy~\cite{borutapy} offers an equivalent Python implementation to the original R implementation, while BorutaSHAP~\cite{borutashap} replaces impurity-based importance with SHAP-value~\cite{lundberg2017shap} based importances. However, none of the available variants address two fundamental limitations of the algorithm. 

First, Boruta is computationally expensive. Each iteration requires training a tree-based ensemble on a dataset augmented with $p$ shadow features, and the algorithm may require hundreds or thousands of iterations to resolve all features. In modern industrial and scientific applications where $p$ is often in the thousands, running Boruta to completion can require days of computation, making the method impractical in time-constrained settings. Second, Boruta provides no convergence (the state in which all features are either confirmed or rejected) guarantees. That is, features can remain \emph{tentative}, neither confirmed nor rejected, forcing practitioners to either increase the iteration budget arbitrarily, enable heuristic early stopping mechanisms or resolve remaining tentatives using post-hoc heuristics such as the rough fix~\cite{kursa2010feature}. In practice, this uncertainty requires careful tuning of hyperparameters that control the maximum number of iterations and early-stopping, creating a tension between computational cost and quality of the solution.

% In this paper, we propose Greedy Boruta, a modification that replaces Boruta's binomial test confirmation criterion with a single-hit rule: a feature is confirmed if it ever exceeds the maximum shadow importance in any iteration. This change is motivated by the fact that, in the all-relevant paradigm, false negatives (missed relevant features) are costlier than false positives (incorrectly included noise features), since false positives can be filtered out by downstream regularization, cross-validation or minimal-optimal selection methods. Despite its simplicity, this modification yields theoretical and computational properties, namely: guaranteed convergence in few iterations, provably equal or higher feature recall (i.e., true positive rate) than Vanilla Boruta, embarrassingly parallel execution, simplification of multiple-testing corrections and elimination of all termination-related hyperparameters. Throughout this paper, we refer to the original algorithm and its equivalent reimplementations collectively as \emph{Vanilla Boruta}, and to our proposed modification as \emph{Greedy Boruta}.

In this paper, we propose Greedy Boruta, a modification that replaces Boruta's binomial-test confirmation criterion with a single-hit rule: a feature is confirmed if it exceeds the maximum shadow importance in at least one iteration. This change is motivated by the fact that, under the all-relevant paradigm, false negatives (missed relevant features) are often more costly than false positives (incorrectly included noise features), since false positives can be filtered by downstream regularization, cross-validation, stability analyses or minimal-optimal feature selection procedures. We further introduce two practical extensions. First, we partition confirmed features into \emph{strongly confirmed} and \emph{weakly confirmed} groups. Strongly confirmed features satisfy the original Boruta confirmation criterion and therefore would also be selected by Vanilla Boruta, whereas weakly confirmed features are selected exclusively by Greedy Boruta and may be viewed as additional candidate variables for further investigation. Second, we show that the number of shadow features can be reduced from $p$ to $r \leq p$, substantially decreasing the dimensionality of the augmented dataset and yielding further computational savings in high-dimensional settings. Despite its simplicity, Greedy Boruta retains several attractive theoretical and computational properties, including (i) guaranteed convergence in a finite number of iterations, (ii) provably equal or higher feature recall than Vanilla Boruta, (iii) embarrassingly parallel execution, (iv) simplified multiple-testing corrections and (v) the elimination of all termination-related hyperparameters. Throughout this paper, we refer to the original algorithm and its equivalent reimplementations collectively as \emph{Vanilla Boruta}, and to our proposed modification as \emph{Greedy Boruta}.


\medskip 
The main contributions of this paper are:
\begin{enumerate}
    \item \textbf{Algorithmic:} A modification to Boruta that guarantees convergence in exactly $K = \lceil \log_2(p/\alpha) \rceil$ iterations\footnote{\label{fn:bonf}Under Bonferroni correction.}, eliminating all termination-related hyperparameters. We further introduce a weak/strong confirmation framework that distinguishes features satisfying the original Boruta confirmation criterion from those selected exclusively by the greedy criterion.
    \item \textbf{Theoretical:} Proofs of (i) a recall guarantee showing that the set of selected features by Vanilla Boruta is contained in the set of features selected by Greedy Boruta (Theorem~\ref{thm:recall}), (ii) a convergence bound on the number of iterations (Theorem~\ref{thm:convergence}), (iii) triviality of false discovery rate (FDR) correction under the greedy criterion (Theorem~\ref{thm:fdr_trivial}), (iv) a false positive rate (FPR) upper bound of $K/(p+1)$ (Theorem~\ref{thm:fpr}) and (v) an analysis of the effect of reducing the number of shadow features from \(p\) to \(r \leq p\) on the computational and statistical properties of Greedy Boruta.
    \item \textbf{Practical:} A reduced-shadow variant of Greedy Boruta that decreases the dimensionality of the augmented dataset from \(2p\) to \(p+r\), yielding additional computational savings while allowing practitioners to explicitly control the trade-off between runtime and false positive rate. The deterministic iteration count and independence across iterations enable embarrassingly parallel execution, implemented in \texttt{greedyboruta} PyPI package.
    \item \textbf{Empirical:} Evaluation on a synthetic benchmark with known ground truth and six real-world high-dimensional datasets, demonstrating 2--66$\times$ speedups with downstream model performance that is statistically indistinguishable from Vanilla Boruta in most configurations.
\end{enumerate}


\section{Background: The Boruta Algorithm}
\label{sec:background}

The idea behind Vanilla Boruta is to construct \emph{shadow features}, permuted copies of features that serve as competitive benchmarks to the non-shadow features. By comparing each non-shadow feature's importance against the maximum shadow importance across multiple randomized trials, Boruta makes statistically principled decisions about which features carry genuine signal and which are not statistically distinguishable from noise.

The algorithm operates iteratively, with each iteration consisting of four steps.
\begin{enumerate}
    \item \textbf{Shadow feature construction}: for each of the $p$ original features, Boruta creates a shadow copy by randomly permuting the feature's values across samples. This permutation destroys any relationship with the target variable while preserving the feature's marginal distribution. The dataset is then extended to include both the $p$ original features and $p$ shadow features. 
    \item \textbf{Importance scoring}: a tree-based ensemble (originally random forests~\cite{breiman2001random}) is trained on this extended dataset, and feature importance scores are extracted for all $2p$ features. 
    \item \textbf{Hit registration}: each original feature whose importance exceeds the maximum shadow importance in that iteration receives a `hit', and these hits are accumulated across iterations. 
    \item \textbf{Statistical testing}: binomial tests determine whether the features are confirmed (relevant), tentative (unknown) or rejected (irrelevant).
\end{enumerate}

The statistical testing framework employs two one-sided binomial tests at each iteration. For a feature with $h$ hits across $K$ iterations, the confirmation test checks whether $\mathbb{P}(X \geq h \mid K, p=0.5) < \alpha$, where $\alpha$ denotes statistical significance and $p = 0.5$ represents the null hypothesis that the feature is no better than random noise (represented by the best shadow). The symmetric rejection test checks $\mathbb{P}(X \leq h \mid K, p=0.5) < \alpha$. Both tests are typically applied with a feature-wise Bonferroni correction to control the family-wise error rate across the $p$ features under test. Features passing neither test remain tentative.

Because there are no convergence guarantees with Vanilla Boruta, when the algorithm terminates with tentative features unresolved, whether by early stopping or reaching a maximum number of iterations, features that remain tentative are resolved using the rough fix heuristic \cite{kursa2010feature}. Under this heuristic, a tentative feature is confirmed if its median importance across iterations exceeds the median of the per-iteration maximum shadow importances, and rejected otherwise. 

\section{The Greedy Boruta Algorithm}
\label{sec:method}

We now present Greedy Boruta, a modification of Vanilla Boruta that achieves dramatic computational speedups while maintaining theoretical guarantees on feature recall. The asymmetry between false negatives (missed relevant features) and false positives (incorrectly confirmed noise features) discussed in Section~\ref{sec:introduction} motivates a relaxation of the confirmation criterion: rather than requiring statistically significant evidence that a feature consistently outperforms random noise, we confirm a feature upon the first indication that it carries signal. This change yields convergence in exactly $K = \lceil \log_2(p/\alpha) \rceil$ iterations (with Bonferroni correction) or $K = \lceil \log_2(1/\alpha) \rceil$ iterations (without), eliminates the need for hyperparameters governing termination, enables embarrassingly parallel execution and provably never misses features that Vanilla Boruta would confirm.

Algorithm~\ref{alg:greedy_boruta} presents the complete Greedy Boruta procedure. Given a dataset with $p$ features and a target variable $y$, the algorithm first determines the total number of iterations $K$, which depends on the number of features and significance level. At each iteration, shadow features are generated by randomly permuting the values of every original feature, thereby preserving marginal distributions while destroying any association with the target. A tree-based ensemble model is then trained on the augmented dataset containing both original and shadow features, and feature importance scores are computed. The maximum importance score among the shadow features defines a noise threshold for that iteration. Any original feature whose importance exceeds this threshold receives a hit. After all iterations are completed, features with at least one hit are confirmed as relevant, whereas features with zero hits are rejected.

The key difference relative to Vanilla Boruta lies in the confirmation step, while the rejection criterion remains effectively unchanged. Vanilla Boruta employs sequential binomial hypothesis tests to determine whether a feature should be confirmed or rejected, whereas Greedy Boruta confirms any feature that achieves at least one hit across all iterations. This modification appears in lines~$14$ and~$15$ of Algorithm~\ref{alg:greedy_boruta}, where the greedy confirmation rule replaces the statistical confirmation test. Furthermore, because a feature with zero hits never exceeds the maximum shadow importance in any iteration, rejection can be determined directly at the end of the procedure without repeatedly performing intermediate statistical tests. Consequently, Greedy Boruta eliminates the sequential testing mechanism used in Vanilla Boruta, resulting in a substantially simpler and more computationally efficient algorithm while preserving the same underlying principle of comparing real features against randomized shadow features.

\begin{algorithm}[t]
\caption{Greedy Boruta Algorithm}
\label{alg:greedy_boruta}
\begin{algorithmic}[1]
\REQUIRE Dataset $X \in \mathbb{R}^{n \times p}$, target $y$, significance level $\alpha$
\ENSURE Set of confirmed features $\mathcal{C}$, rejected features $\mathcal{R}$
\STATE Compute $K = \lceil \log_2(p/\alpha) \rceil$ \COMMENT{Total iterations}
\STATE Initialize $\text{hits}[j] = 0$ for all features $j = 1, \ldots, p$
\FOR{$i = 1$ to $K$}
    \STATE Create shadow features by permuting each column of $X$
    \STATE Train tree ensemble on $[X, X_{\text{shadow}}]$ with target $y$
    \STATE Compute importance scores for all features
    \STATE $m \gets \max(\text{importance of shadow features})$
    \FOR{each original feature $j$}
        \IF{$\text{importance}[j] > m$}
            \STATE $\text{hits}[j] \gets \text{hits}[j] + 1$
        \ENDIF
    \ENDFOR
\ENDFOR
\STATE $\mathcal{C} \gets \{j : \text{hits}[j] \geq 1\}$ \COMMENT{Greedy confirmation}
\STATE $\mathcal{R} \gets \{j : \text{hits}[j] = 0\}$ \COMMENT{All others rejected ($p-value < \alpha/p$)} 
\RETURN $\mathcal{C}, \mathcal{R}$
\end{algorithmic}
\end{algorithm}


\subsection{Weak and strong confirmations}

A potential limitation of the Greedy Boruta confirmation rule is that all features with at least one hit are treated equally, regardless of the total number of hits accumulated across iterations. In particular, a feature that beats the shadow threshold once receives the same confirmation status as a feature that consistently outperforms the shadow threshold. To provide a more nuanced interpretation of the selected features while preserving the feature recall guarantee established in Theorem~\ref{thm:recall}, which is presented in Section \ref{sec:method:theory}, we propose partitioning the confirmed features into \emph{strongly} and \emph{weakly} confirmed groups.

Let \(H_j\) denote the number of hits accumulated by feature \(j\) after \(K\) iterations. Recall that Greedy Boruta confirms a feature whenever \(H_j \geq 1\). Define the set of confirmed features as
\[
\mathcal{C}_G
=
\{j : H_j \geq 1\}.
\]
We further define the set of \emph{strongly confirmed} features as
\[
\mathcal{C}_{\mathrm{strong}}
=
\left\{
j :
\mathbb{P}\!\left(
X \geq H_j
\right)
\leq \alpha,
\quad
X \sim \mathrm{Binomial}(K,0.5)
\right\}.
\]
That is, a feature is strongly confirmed if it satisfies the original Boruta confirmation criterion. Equivalently, these are the features that would be confirmed by Vanilla Boruta after \(K\) iterations. 

The remaining confirmed features are classified as \emph{weakly confirmed} and are part of the set
\[
\mathcal{C}_{\mathrm{weak}}
=
\left\{
j :
H_j \geq 1
\ \text{and}\
\mathbb{P}\!\left(
X \geq H_j
\right)
>
\alpha,
\quad
X \sim \mathrm{Binomial}(K,0.5)
\right\}.
\]
The distinction between weak and strong confirmations provides a natural ranking of the selected features. Strongly confirmed features have accumulated sufficient evidence to satisfy the original Boruta significance criterion and therefore correspond to features that would also be selected by Vanilla Boruta. In contrast, weakly confirmed features are selected exclusively by Greedy Boruta. These features have demonstrated some evidence of relevance by outperforming the shadow threshold at least once, but the evidence is not strong enough to achieve statistical significance under the Boruta binomial test.

From a practical perspective, strongly confirmed features may be interpreted as high-confidence discoveries, whereas weakly confirmed features may warrant additional investigation through domain expertise, stability analyses or downstream predictive validation. This distinction allows Greedy Boruta to retain its computational advantages and feature recall guarantee while providing users with a measure of confidence among the selected features.

\subsection{Reducing the Number of Shadow Features}

By relaxing the confirmation criterion for relevant features while preserving the statistical rigour of the rejection criterion, Greedy Boruta achieves substantial reductions in computational time across the datasets considered in Section~\ref{sec:experiments}. A key advantage of Greedy Boruta is that the total number of iterations can be determined analytically from the significance level~\(\alpha\) and the number of features $p$, unlike Vanilla Boruta, whose convergence behaviour depends on sequential hypothesis testing and heuristic stopping rules. In this section, we further investigate how the computational cost of Greedy Boruta can be reduced by decreasing the number of shadow features used to define the random-noise threshold at each iteration. To disWe refer to this variant

% To quantify the effect of using only $r$ shadow features, we first establish the following result, which the proof is given in Appendix~\ref{app:proofs}.

% \begin{proposition}
% \label{prop:mr_mp}
% Let \(S_1,\ldots,S_p\) be i.i.d.\ continuous random variables denoting shadow importances with cumulative distribution function \(F\). Define
% \[
% M_r = \max(S_1,\ldots,S_r),
% \qquad
% M_p = \max(S_1,\ldots,S_p),
% \]
% for \(r \leq p\). Then
% \[
% \mathbb{P}(M_r = M_p) = \frac{r}{p}.
% \]
% \end{proposition}
% \begin{remark}
% The result above quantifies how often the reduced-shadow threshold coincides exactly with the threshold that would have been obtained using all $p$ shadows. In practice, the reduced-shadow threshold agrees with the one from vanilla Boruta with probability $r/p$.
% \end{remark}

% \begin{remark}
% The result also follows immediately from symmetry. Among the \(p\) shadow importances, the global maximum $M_p$ is equally likely to occur at any index. Therefore, the probability that the global maximum belongs to the subset of size \(r\) is \(r/p\).
% \end{remark}


% Recall that both Vanilla Boruta and Greedy Boruta augment the original dataset containing \(p\) features with \(p\) shadow features obtained through random permutations of the original covariates. Consequently, as shown in Algorithm~\ref{alg:greedy_boruta}, the base estimator is trained on a design matrix of dimension \(n \times 2p\). However, when \(p\) is large, constructing all \(p\) shadow features may become computationally expensive. To address this, we consider using only \(r \leq p\) shadow features instead of the full set of \(p\) shadows. In practice, at each iteration, we generate only $r$ shadow features obtained by randomly permuting $r$ selected covariates. The random-noise threshold is defined as the maximum importance among the generated shadow features. This modification is particularly attractive in high-dimensional settings, where the cost of fitting the base estimator dominates the overall runtime. In addition, this may be important when shadow feature creation and storage is most costly, for example, when the augmented dataset does not fully fit in RAM and overhead regarding materializing to disk is added.

Recall that both Vanilla Boruta and Greedy Boruta augment the original dataset containing \(p\) features with \(p\) shadow features obtained through random permutations of the original covariates. Consequently, as shown in Algorithm~\ref{alg:greedy_boruta}, the base estimator is trained on a design matrix of dimension \(n \times 2p\). However, when \(p\) is large, constructing all \(p\) shadow features may become computationally expensive. To address this, we consider using only \(r \leq p\) shadow features instead of the full set of \(p\) shadows. In practice, at each iteration, we generate only \(r\) shadow features obtained by randomly permuting \(r\) selected covariates. The random-noise threshold is then defined as the maximum importance among the generated shadow features.

This modification is particularly attractive in high-dimensional settings, where the cost of fitting the base estimator dominates the overall runtime. Moreover, reducing the number of shadow features can substantially decrease memory and storage requirements. This is especially relevant in distributed computing environments for big data, such as Apache Spark \cite{zaharia2016spark}, where the augmented dataset may not fit entirely in memory. In such settings, shadow features either need to be represented explicitly in the execution plan or materialized as intermediate data on disk, both of which introduce additional computational and storage overhead. Using fewer shadow features alleviates these costs while preserving the statistical framework of the method.

% As discussed in Section~\ref{sec:method:fpr}, under Assumption~\ref{asm:exchangeability}, the importance of an irrelevant feature and the \(p\) shadow importances form \(p+1\) exchangeable random variables. Consequently, when all \(p\) shadow features are used, each variable is equally likely to attain the maximum importance score, implying that the per-iteration hit probability of an irrelevant feature is
% \[
% \theta_0 = \frac{1}{p+1}.
% \]


% Assume $S_1,\dots,S_p$ are i.i.d. continuous random variables denoting shadow feature importances with cumulative distribution function \(F\), and define
% $$
% M_p = \max_{1 \leq j \leq p} S_j,
% \qquad
% M_r = \max_{1 \leq j \leq r} S_j.
% $$
% Conditioning on the observed shadow importances,
% \[
% \mathbb{P}(M_r=M_p \mid S_1,\ldots,S_r)
% =
% F(M_r)^{p-r},
% \]
% since all remaining \(p-r\) shadow importances must fall below \(M_r\).

% \begin{enumerate}
%     \item For any \textit{observed importance} \(i \leq r\),
% $$
% \mathbb{P}(S_i = M_p \mid S_1,\dots,S_r)
% =
% \mathbf{1}\{S_i = M_r\}\,F(S_i)^{p-r}.
% $$
% That is, if \(S_i \neq M_r\), then \(S_i\) cannot be the global maximum. Otherwise, if \(S_i = M_r\), then all remaining \(p-r\) variables must satisfy $S_{r+1},\dots,S_p < S_i$, which occurs with probability \(F(S_i)^{p-r}\). \\


% \item For any \textit{unobserved importance} \(i > r\),
% $$
% \mathbb{P}(S_i = M_p \mid S_1,\dots,S_r)
% =
% \int_{M_r}^{\infty} F(x)^{p-r-1}\,dF(x).
% $$
% Using the substitution \(u = F(x)\),
% \begin{align}
% \mathbb{P}(S_i = M_p \mid S_1,\dots,S_r) & = \int_{F(M_r)}^{1} u^{p-r-1}\,du \nonumber \\
% & = \frac{1 - F(M_r)^{p-r}}{p-r} \nonumber.
% \end{align}
% Summing over $p-r$ unobserved shadows gives
% $$
% 1 - F(M_r)^{p-r},
% $$
% which is the probability that some unseen shadow exceeds $M_r$. Therefore, $F(M_r)^{p-r}$ is the probability that the maximum among the sampled shadows is also the global shadow maximum (i.e., $M_r = M_p)$.
% \end{enumerate}
Let \(\theta_0^\star\) denote the per-iteration hit probability of an irrelevant feature when only \(r\) shadow features are used. Under Assumption~\ref{asm:exchangeability}, which is formally introduced in Section \ref{sec:method:theory}, the importance \(I_j\) of an irrelevant feature and the \(r\) shadow importances form \(r+1\) exchangeable random variables. Consequently,
\[
\theta_0^\star = \frac{1}{r+1},
\]
which recovers the result presented in Section \ref{sec:method:fpr} when \(r=p\), where \(\theta_0 = 1/(p+1)\). Since the iterations are independent, \(H_j \sim \mathrm{Binomial}(K,\theta_0^\star)\), and under the greedy criterion with \(r\) shadow features, a false positive occurs whenever an irrelevant feature receives at least one hit. Applying the same argument as in the proof of Theorem~\ref{thm:fpr}, the false positive probability is
\begin{align}
\label{eqn:fpr_upper_bound}
\mathbb{P}(H_j \geq 1)
&=
1-(1-\theta_0^\star)^K \nonumber\\
&=
1-\left(1-\frac{1}{r+1}\right)^K \nonumber\\
&\leq \frac{K}{r+1}. 
\end{align}
This result highlights a direct trade-off between computational cost and false positive control. Reducing the number of shadow features decreases the dimensionality of the augmented design matrix from \(2p\) to \(p+r\), potentially yielding substantial computational savings. However, the null hit probability increases from \(1/(p+1)\) to \(1/(r+1)\), resulting in a larger FPR.

Figure~\ref{fig:relative_fpr_shadows} illustrates this trade-off for different combinations of number of features \(p\), significance level \(\alpha\) and fraction of shadow features \(r/p\). The impact of reducing the number of shadow features depends strongly on the dimensionality of the problem. For small values of \(p\), the FPR is already relatively large when all shadow features are used, so further reductions in \(r\) may lead to an undesirable increase in spurious confirmations. For example, when \(p=500\) and \(\alpha=0.05\), the per-feature FPR is approximately \(2.8\%\) when \(r=p\) and increases $4.6$ times to \(13\%\) when only \(20\%\) of the shadow features are used. In contrast, for moderate and large values of \(p\), the FPR remains small even when only a fraction of the shadow features is used. For instance, when \(p=10{,}000\) and \(\alpha=0.05\), the FPR is approximately \(0.18\%\) when \(r=p\) and remains below \(1\%\) when only \(20\%\) of the shadow features are used. In practice, an FPR of \(\gamma\) indicates that a truly irrelevant feature is incorrectly confirmed with probability \(\gamma\). Therefore, if there are \(m_0\) irrelevant features, the expected number of false positives is \(m_0\gamma\).

\begin{figure}
    \centering
    \includegraphics[width=\linewidth]{relative_FPR_with_shadows.png}
    \caption{Per-feature false positive rate (FPR) relative to the minimum FPR across all $\alpha$ at $100\%$ of Greedy Boruta under Bonferroni correction using different number of features, fraction of shadow features and significance level.}
    \label{fig:relative_fpr_shadows}    
\end{figure}

A practical question is how to choose the number of shadow features \(r\). Given the upper bound on the FPR in Equation \eqref{eqn:fpr_upper_bound}, one could, in principle, select \(r\) so that the upper bound does not exceed a significance level \(\alpha\). Solving
\[
\frac{K}{r+1}\leq \alpha
\]
yields the condition $r \geq K/\alpha-1$. Therefore, a conservative choice is
\[
r=
\min\!\left\{
p,\,
\left\lceil \frac{K}{\alpha}-1 \right\rceil
\right\}.
\]
This guarantees that the theoretical upper bound on the per-feature FPR is at most \(\alpha\). However, because \(K=\lceil \log_2(p/\alpha)\rceil\), this rule can recommend \(r=p\) for practical values of \(p\) and \(\alpha\), providing little computational benefit. As illustrated in Figure~\ref{fig:relative_fpr_shadows}, however, substantial reductions in the number of shadow features can often be achieved in medium to high-dimensional settings while maintaining an acceptably low FPR.

Additional results are provided in Appendix~\ref{app:add_results_reducing_shadows}. In particular, Figure~\ref{fig:fpr_shadows} and Tables~\ref{tab:fpr_shadows} and~\ref{tab:fpr_shadows_exact} report the theoretical FPR and their corresponding upper bounds for different combinations of the number of features \(p\), the fraction of shadow features \(r/p\) and the significance level \(\alpha\).

\subsection{Theoretical Properties}
\label{sec:method:theory}

\paragraph{Maintained Recall.}
As we prove in Theorem~\ref{thm:recall}, because confirmation becomes easier (only one hit required vs. statistical significance), Greedy Boruta can never have a lower true positive rate (i.e., feature recall) than Vanilla Boruta. Any feature confirmed by Vanilla Boruta must have accumulated sufficient hits to achieve statistical significance, which necessarily means it has at least 1 hit. Therefore, any feature confirmed by the vanilla approach is also confirmed by greedy one. It is important to note that this assumes both the vanilla and greedy algorithms are either using or not using bonferroni correction.

\paragraph{Guaranteed Convergence.}
As we prove in Theorem~\ref{thm:convergence}, the greedy criterion allows computing \emph{a priori} the exact number of iterations required for convergence, solving the uncertain termination of Vanilla Boruta. Notably, the minimum number of iterations for Vanilla Boruta to make its first confirmation or rejection in a given task equals the total iteration count at which Greedy Boruta completes for that same task.

\paragraph{Hyperparameter Simplification.}
% The guaranteed convergence eliminates the \texttt{max\_iter}, \texttt{early\_stopping}, and \texttt{n\_iter\_no\_change} parameters from \texttt{boruta\_py} (Corollary~\ref{cor:hyperparams}).
A direct consequence of the guaranteed convergence is a hyperparameter simplification in the implementation of Greedy Boruta. Namely, the \texttt{max\_iter} (maximum number of Boruta iterations), \texttt{early\_stopping} (heuristic stopping rule), and \texttt{n\_iter\_no\_change} (how many consecutive iterations with no change in feature status are tolerated) parameters present in the \texttt{boruta\_py} package are not required.

\paragraph{Parallelization.}
\label{sec:method:parallel}
Because the number of iterations is fixed at $K$ and no features are removed mid-execution, all iterations operate on an identical feature set with independent random seeds. The hit accumulation is commutative and associative, so all $K$ iterations can be executed in parallel with results aggregated at the final classification step. In contrast, Vanilla Boruta is inherently sequential: features may be dropped mid-run, early stopping monitors consecutive iterations and statistical tests accumulate evidence over time. We implement this parallelization in \texttt{greedyboruta} using \texttt{joblib}. The observed speedup depends on the base estimator's internal parallelism: when tree training already saturates all CPU cores (e.g., \texttt{ExtraTreesClassifier} with \texttt{n\_jobs=-1}), outer iteration-level parallelism yields limited additional benefit, and the dominant speedup comes from the reduced iteration count.

We now state the theoretical guarantees of Greedy Boruta. Full proofs are provided in Appendix~\ref{app:proofs}.

\subsubsection{Feature Recall Guarantee}

\begin{theorem}[Feature Recall Guarantee]
\label{thm:recall}
Let $\mathcal{C}_V$ and $\mathcal{C}_G$ denote the sets of features confirmed by Vanilla Boruta and Greedy Boruta, respectively, when run with the same paramaters, dataset, model and both with or without the Bonferroni correction. Then, 
$$\mathcal{C}_V \subseteq \mathcal{C}_G.$$ 
The inclusion holds regardless of whether Vanilla Boruta's confirmations arise from the binomial test or from the rough fix heuristic applied to tentative features at termination. 
\end{theorem}

The key insight is that both confirmation mechanisms in Vanilla Boruta (the binomial test and the rough fix) require a feature to have at least one hit. For the binomial test, zero hits yield a p-value of 1, which cannot be rejected. For the rough fix, a feature with zero hits has its empirical importance distribution stochastically dominated \cite{levy1990stochastic} by the maximum shadow importance, so its median cannot exceed the shadow median. Since the greedy criterion confirms any feature with at least one hit, every feature confirmed by Vanilla Boruta is also confirmed by Greedy Boruta.

\subsubsection{Convergence Guarantee}

\begin{theorem}[Convergence Guarantee]
\label{thm:convergence}
For a dataset with $p$ features, the Greedy Boruta algorithm terminates with no tentative features in exactly $K = \lceil \log_2(p/\alpha) \rceil$ iterations when using Bonferroni correction for multiple comparisons.
\end{theorem}

At iteration $K$, every feature with at least 1 hit is confirmed by the greedy criterion. Every remaining feature has zero hits and thus a rejection p-value of $(0.5)^K$. Under Bonferroni correction, rejection requires $(0.5)^K < \alpha/p$, which holds by construction of $K = \lceil \log_2(p/\alpha) \rceil$. Therefore, all features are classified and the algorithm terminates. It is worth noting that the number of iterations decomposes as $K = \lceil \log_2(p) + \log_2(1/\alpha) \rceil$: halving $\alpha$ or doubling $p$ adds exactly one iteration. Table~\ref{tab:iterations_alpha} shows $K$ for representative values.

\begin{table}[H]
\centering
\caption{Number of iterations $K = \lceil \log_2(p/\alpha) \rceil$ required for convergence when using Bonferroni correction for multiple comparisons.}
\label{tab:iterations_alpha}
\begin{tabular}{c|ccccc}
\toprule
& \multicolumn{5}{c}{Number of features $p$} \\
$\alpha$ & 50 & 100 & 500 & 1000 & 5000 \\
\midrule
0.100 & 9 & 10 & 13 & 14 & 16 \\
0.050 & 10 & 11 & 14 & 15 & 17 \\
0.010 & 13 & 14 & 16 & 17 & 19 \\
0.001 & 16 & 17 & 19 & 20 & 23 \\
$10^{-4}$ & 19 & 20 & 23 & 24 & 26 \\
$10^{-5}$ & 23 & 24 & 26 & 27 & 29 \\
\bottomrule
\end{tabular}
\end{table}

\noindent Without Bonferroni correction (analogous to \texttt{mcAdj=FALSE} in the R implementation), the iteration count reduces to $K = \lceil \log_2(1/\alpha) \rceil$, and the number of iterations is not dependent on the number of features anymore. In Sections \ref{sec:experiments:setup} and \ref{sec:experiments:realworld}, we use the Bonferroni correction to obtain $K$. 

\subsubsection{Triviality of FDR Correction in Greedy Boruta}

\begin{theorem}[FDR Correction Triviality]
\label{thm:fdr_trivial}
At iteration $K$ in Greedy Boruta, all tentative features (those with 0 hits) have identical p-values equal to $(0.5)^K$. Consequently, the Benjamini-Hochberg FDR correction \cite{benjamini1995controlling} either rejects all tentative features or none, making the correction equivalent to an uncorrected test.
\end{theorem}

This follows directly from the greedy confirmation criterion: all features with at least 1 hit are already confirmed, so the only features entering the rejection test have exactly zero hits and identical p-values. The Benjamini-Hochberg procedure \cite{benjamini1995controlling} applied to a set of identical p-values reduces to a single threshold comparison. This contrasts with Vanilla Boruta, where tentative features accumulate different hit counts and thus have a distribution of p-values, making FDR correction non-trivial.

\subsection{Multiple Testing Corrections}
\label{sec:method:corrections}

% Greedy Boruta simplifies the multiple testing correction done in Boruta. Because confirmation requires no statistical test and rejection occurs only once at iteration $K$, trial-wise (temporal) correction is unnecessary. Furthermore, by Theorem~\ref{thm:fdr_trivial}, all features entering the rejection test have identical p-values, making FDR correction equivalent to an uncorrected test. The only meaningful correction is feature-wise Bonferroni for rejection.

Table~\ref{tab:correction_mapping} summarizes how the correction settings in \texttt{greedyboruta} relate to those available in the most widely used Boruta implementations. In Vanilla Boruta, multiple-testing correction is required both across features and across iterations because confirmation and rejection decisions are repeatedly revisited throughout the execution of the algorithm. In contrast, Greedy Boruta performs a single rejection test at the final iteration and does not rely on hypothesis testing for confirmation. Consequently, trial-wise correction becomes unnecessary. Furthermore, Theorem~\ref{thm:fdr_trivial} shows that all features entering the rejection stage have identical p-values, causing FDR-based procedures to degenerate to the uncorrected test. As a result, the only non-trivial correction remaining in Greedy Boruta is the feature-wise Bonferroni adjustment, controlled through the \texttt{bonferroni} parameter.

\begin{table}[t]
\centering
\caption{Mapping between correction settings across implementations.}
\label{tab:correction_mapping}
\begin{tabular}{lll}
\toprule
\texttt{greedyboruta} & R Boruta & \texttt{boruta\_py} \\
\midrule
\texttt{bonferroni=True}  & \texttt{mcAdj=TRUE}  & \texttt{two\_step=False} \\
\texttt{bonferroni=False} & \texttt{mcAdj=FALSE} & \texttt{two\_step=True}\footnotemark \\
\bottomrule
\end{tabular}
\end{table}

\footnotetext{The \texttt{two\_step=True} mode applies both FDR and trial-wise Bonferroni correction; both degenerate in Greedy Boruta (Theorem~\ref{thm:fdr_trivial}), reducing to the uncorrected case.}

\subsection{False Positive Rate Analysis}
\label{sec:method:fpr}

The feature recall and convergence guarantees come at a cost: by relaxing the confirmation criterion, Greedy Boruta may confirm features that Vanilla Boruta would leave tentative or reject. We quantify this trade-off under a standard assumption of the Boruta framework.

\begin{assumption}[Shadow Exchangeability]
\label{asm:exchangeability}
An irrelevant feature is exchangeable with the shadow features with respect to the importance measure. That is, the importance of an irrelevant feature and the importances of the $p$ shadow features are identically distributed and symmetric under permutation within each iteration.
\end{assumption}

Violations of this assumption (e.g., importance measure biases toward high-cardinality features~\cite{strobl2007bias}) affect both Vanilla Boruta and Greedy Boruta equally, since both use the same shadow mechanism. The analysis below therefore isolates the effect of the modified decision rule.

Under the shadow exchangeability assumption, the importance of an irrelevant feature and the $p$ shadow importances form $p+1$ exchangeable random variables, so each is equally likely to be the maximum. The per-iteration hit probability for an irrelevant feature is therefore $\theta_0 = 1/(p+1)$.

\begin{theorem}[FPR of Greedy Boruta]
\label{thm:fpr}
Under Assumption~\ref{asm:exchangeability}, the probability that Greedy Boruta falsely confirms an irrelevant feature is
\begin{equation}
\label{eq:fpr_exact}
    \mathbb{P}(\text{false positive}) = 1 - \left(\frac{p}{p+1}\right)^{K},
\end{equation}
where $K = \lceil \log_2(p/\alpha) \rceil$. This quantity is bounded above by
\begin{equation}
\label{eq:fpr_bound}
    \mathbb{P}(\text{false positive}) \leq \frac{K}{p + 1} = \frac{\lceil \log_2(p/\alpha) \rceil}{p + 1}.
\end{equation}
If there are $m_0$ independent irrelevant features, the expected false positive count satisfies $\mathbb{E}[\text{FP}] \leq m_0 \cdot K / (p+1)$.
\end{theorem}

The proof, given in Appendix~\ref{app:proofs}, follows from the binomial distribution of hits and Bernoulli's inequality. The bound has a favorable structure: the numerator grows logarithmically while the denominator grows linearly, so the FPR \emph{decreases} as $p$ increases. Table~\ref{tab:fpr} presents the bound for representative configurations.

\begin{table}[t]
\caption{Upper bound on per-feature false positive rate (FPR) of Greedy Boruta under Bonferroni correction ($K = \lceil \log_2(p/\alpha) \rceil$). Each cell shows the value obtained with the FPR bound equation.}
\label{tab:fpr}
\centering
\begin{tabular}{c|cc|cc|cc|cc}
\toprule
& \multicolumn{2}{c|}{$p = 100$}
& \multicolumn{2}{c|}{$p = 500$}
& \multicolumn{2}{c|}{$p = 1{,}000$}
& \multicolumn{2}{c}{$p = 10{,}000$} \\
$\alpha$
& $K$ & FPR
& $K$ & FPR
& $K$ & FPR
& $K$ & FPR \\
\midrule
$0.05$    
& 11 & $10.89\%$
& 14 & $2.79\%$
& 15 & $1.50\%$
& 18 & $0.18\%$ \\

$0.01$    
& 14 & $13.86\%$
& 16 & $3.19\%$
& 17 & $1.70\%$
& 20 & $0.20\%$ \\

$0.001$   
& 17 & $16.83\%$
& 19 & $3.79\%$
& 20 & $2.00\%$
& 24 & $0.24\%$ \\

$0.0001$  
& 20 & $19.80\%$
& 23 & $4.59\%$
& 24 & $2.40\%$
& 27 & $0.27\%$ \\

$0.00001$ 
& 24 & $23.76\%$
& 26 & $5.19\%$
& 27 & $2.70\%$
& 30 & $0.30\%$ \\
\bottomrule
\end{tabular}
\end{table}

For moderate to large $p$ (the typical Boruta use case), the bound is comfortably small: below $5\%$ for $p \geq 500$ across all $\alpha$ values shown. For small $p$ (e.g., $p = 100$), the bound can exceed $10\%$, but in such settings the computational cost of Vanilla Boruta is typically manageable.


\section{Experimental Evaluation}
\label{sec:experiments}
 
We evaluate Greedy Boruta on two complementary fronts. First, using a synthetic dataset with a known ground truth, we measure selection quality directly against the true relevant feature set across a sweep of significance levels~$\alpha$ and analyse how Vanilla Boruta's quality evolves over iterations relative to the point at which Greedy Boruta terminates by design. Second, we validate on a suite of real-world high-dimensional benchmarks where ground truth is
unavailable, using downstream model performance as a proxy for selection quality. We point out that all experiments in this section use the original Boruta setting with $p$ shadow features. Results for the reduced-shadow variant ($r \leq p$) of Greedy Boruta, which explicitly trades false positive control for additional computational savings, are presented in Appendix~\ref{app:add_results_reducing_shadows}.

We compare Greedy Boruta primarily against the R implementation's default configuration (\texttt{mcAdj=TRUE}, applying feature-wise Bonferroni correction). In practive, since the comparison is made with the \texttt{boruta\_py} package, we use \texttt{two\_step=False} to match the R defaults. All experiments use $R=10$ independent random seeds; results are reported as mean~$\pm$~standard deviation (std) unless stated otherwise.
 
% -----------------------------------------------------------------------------
\subsection{Synthetic Dataset Experiments}
\label{sec:experiments:setup}
 
We construct a controlled synthetic benchmark using \texttt{sklearn.datasets.make\_classification} with the following parameters: $n=1{,}000$ samples, $p=500$ features (\texttt{n\_informative = 5}, \texttt{n\_redundant = 100}, \texttt{n\_repeated = 0}, remaining $m_0 = 395$ are pure noise), one cluster per class (\texttt{n\_clusters\_per\_class = 1}), a low class separation (\texttt{class\_sep = 0.0001}) to prevent trivial discrimination, and a 30\% label noise rate (\texttt{flip\_y = 0.3}). All features are standardized to zero mean and unit variance prior to selection. Because the Boruta algorithm is designed for the all-relevant paradigm, we consider both the 5 informative and the 100 redundant features as relevant. We call this set of 105 relevant features the \emph{Broad}~GT (ground truth). The significance level is swept over $\alpha \in \{10^{-5},\,10^{-4},\,10^{-3},\,10^{-2},\,10^{-1}\}$. The base estimator for both methods is \texttt{ExtraTreesClassifier}~\cite{geurts2006extremely,pedregosa2011scikit}(\texttt{max\_depth=5}, \texttt{n\_estimators=`auto'}, \texttt{n\_jobs=-1}). Vanilla Boruta is configured with \texttt{max\_iter=512}, \texttt{early\_stopping=True} (\texttt{n\_iter\_no\_change=128}), and \texttt{two\_step=False}. Greedy Boruta uses Bonferroni correction (\texttt{bonferroni=True}) and \texttt{n\_jobs=4} for outer iteration-level parallelism. All experiments were run on an Intel Core~i7-14700K (3.40~GHz, 20 cores: 8P+12E), and with 64~GB RAM.

\paragraph{Parallelism and speedup decomposition.}
Both methods use the same base estimator configuration with \texttt{n\_jobs=-1}. Vanilla Boruta executes its iterations sequentially; Greedy Boruta additionally distributes iterations across 4 workers, parallelizing iterations as well as tree construction. Consequently, the observed speedup possibly reflects two factors: (i)~fewer iterations ($K = \lceil \log_2(p/\alpha) \rceil$ vs.\ hundreds) and (ii)~outer iteration-level parallelism. However, because each base estimator training call already saturates most of the available cores, effect (i) dominates (ii) in terms of speedup in most cases. Realistically, in non-cluster environments, the marginal benefit of outer parallelism is sub-linear and substantially smaller than the iteration-count reduction.

\paragraph{Performance Metrics.}
We evaluate against Broad~GT using four feature-level binary classification metrics: recall ($\text{TP}/(\text{TP}+\text{FN})$), precision ($\text{TP}/(\text{TP}+\text{FP})$), specificity ($\text{TN}/(\text{TN}+\text{FP})$) and F1 score ($\text{2*TP}/(\text{2*TP}+\text{FP}+\text{FN})$), where a \emph{true positive} (TP) is a relevant feature that was confirmed, a \emph{false positive} (FP) is a noise feature incorrectly confirmed, a \emph{false negative} (FN) is a relevant feature missed and a \emph{true negative} (TN) is a noise feature correctly rejected. Runtime comparisons use wall-clock time in seconds and speedup is defined as the ratio of each Vanilla Boruta run's time to the Greedy Boruta time .
 
\begin{figure}
    \centering
    \includegraphics[width=\textwidth, height=0.8\textheight, keepaspectratio]{alpha_effect_boxplots.png}
    \caption{Final-state ground-truth metrics across significance levels ($\alpha$) for Vanilla
    Boruta and Greedy Boruta on the synthetic dataset. Individual runs are shown as semi-transparent dots; markers with caps show mean $\pm$ SE.}
    \label{fig:final_state_metrics}
\end{figure}
 
Figure~\ref{fig:final_state_metrics} shows the final-state metric distributions. Greedy Boruta achieves equal or higher feature recall than Vanilla Boruta at every $\alpha$ level, in line with Theorem~\ref{thm:recall}. The cost is a modest reduction in precision and specificity: across individual runs, Greedy Boruta incorrectly confirms between 4 and 11 noise features (out of 395), a small fraction of the noise space. Recall is statistically greater for Greedy Boruta (one-sided Welch $t$-test, $p < 0.01$) and precision and specificity are statistically smaller at all $\alpha$ levels, as expected. No statistically significant difference in F1 was detected between the two methods at any $\alpha$ level (two-sided Welch $t$-test, $p > 0.01$), consistent with the competing effects of higher recall and lower precision offsetting each other. Mean speedup ratios (Figure~\ref{fig:final_state_metrics}, bottom panel) fall within the range of 20--35$\times$ across all $\alpha$ values.
 
\subsection{Real-World Benchmarks}
\label{sec:experiments:realworld}
 
To validate Greedy Boruta beyond controlled synthetic settings, we evaluate it on six publicly available high-dimensional benchmark datasets spanning both classification and regression tasks (Table~\ref{tab:realworld_datasets}). Unlike the synthetic setting described in Section~\ref{sec:experiments}, the ground-truth relevant feature sets are unknown for these datasets, preventing direct evaluation of feature-selection accuracy. We therefore assess selection quality indirectly through downstream predictive performance. Both Greedy Boruta and Vanilla Boruta are executed independently using the same base estimator (\texttt{ExtraTreesClassifier} or \texttt{ExtraTreesRegressor}), parallelism configuration and hardware environment. After each method selects its subset of relevant features, a downstream Light Gradient Boosting Machine (LightGBM) model~\cite{ke2017lightgbm} is trained using only the selected features. Performance is then evaluated using 5-fold cross-validation and measured by weighted F1 score (classification) and $R^2$ (regression).

%Each (dataset, $\alpha$, method) configuration is repeated over $R = 10$ independent random seeds; results are reported as mean~$\pm$~SE.
 
\begin{table}
\centering
\caption{Real-world benchmark datasets.}
\label{tab:realworld_datasets}
\begin{tabular}{llrrll}
\toprule
Dataset    & Domain          & $n$    & $p$      & Task           & Source \\
\midrule
Madelon    & Artificial      & 2{,}000  & 500      & Classification & \cite{guyon2004result} \\
Gisette    & Digit recog.    & 6{,}000  & 5{,}000  & Classification & \cite{guyon2004result} \\
Arcene     & Mass spectrom.  & 100      & 10{,}000 & Classification & \cite{guyon2004result} \\
Dexter     & Text classif.   & 300      & 20{,}000 & Classification & \cite{guyon2004result} \\
Isolet     & Speech recog.   & 7{,}797  & 617      & Classification & \cite{uci2017} \\
Riboflavin & Gene expression & 71       & 4{,}088  & Regression     & \cite{buhlmann2014hdi} \\
\bottomrule
\end{tabular}
\end{table}

The significance level is swept over $\alpha \in \{10^{-1},\, 10^{-3},\, 10^{-5}\}$, a reduced set, since  these experiments are much more computationally intensive. All other experimental settings are identical to those described in Section~\ref{sec:experiments:setup}.
 
\subsubsection{Downstream Performance}
\label{sec:experiments:realworld:performance}
 
Table~\ref{tab:realworld_results} presents the downstream performance, runtime and speedup for each dataset and significance level. To assess whether the observed performance differences are statistically meaningful, we conduct one-sided Welch $t$-tests at significance level $0.01$: one testing whether Greedy Boruta is significantly better, and the converse. The ``Sig.'' column reports the outcome: $\approx$ indicates no statistically significant difference (tie), $\blacktriangleright$ indicates Vanilla is significantly better, and $\blacktriangleleft$ indicates Greedy is significantly better.

\begin{table}
\caption{Real-world benchmark results (mean $\pm$ SE over $R=10$ seeds).
\emph{Score} is test-set F1 (classification) or $R^2$ (regression).
\emph{Sig.}:\ directional one-sided Welch $t$-test at level $0.01$;
$\blacktriangleright$~= Vanilla significantly better,
$\blacktriangleleft$~= Greedy significantly better,
$\approx$~= no significant difference.
\emph{Speedup} = Vanilla time\,/\,Greedy time.}
\label{tab:realworld_results}
\centering
\small
\begin{tabular}{ll cc cc c r}
\toprule
& & \multicolumn{2}{c}{Greedy Boruta}
& \multicolumn{2}{c}{Vanilla Boruta} & & \\
\cmidrule(lr){3-4} \cmidrule(lr){5-6}
Dataset & $\alpha$ & Score & Time (s)
                   & Score & Time (s)
                   & Sig. & Speedup \\
\midrule
\multirow{3}{*}{Madelon}
  & $10^{-1}$  & $.856 \pm .002$ & 2.4
               & $.862 \pm .002$ & 34.7
               & $\blacktriangleright$ & $14.6\times$ \\
  & $10^{-3}$  & $.856 \pm .002$ & 3.2
               & $.866 \pm .001$ & 34.8
               & $\blacktriangleright$ & $10.9\times$ \\
  & $10^{-5}$  & $.855 \pm .002$ & 4.1
               & $.866 \pm .001$ & 39.1
               & $\blacktriangleright$ & $9.5\times$ \\
\midrule
\multirow{3}{*}{Arcene}
  & $10^{-1}$  & $.868 \pm .005$ & 11.2
               & $.878 \pm .004$ & 184.3
               & $\approx$ & $16.4\times$ \\
  & $10^{-3}$  & $.875 \pm .005$ & 14.0
               & $.879 \pm .003$ & 203.3
               & $\approx$ & $14.6\times$ \\
  & $10^{-5}$  & $.875 \pm .006$ & 19.9
               & $.878 \pm .005$ & 215.7
               & $\approx$ & $10.9\times$ \\
\midrule
\multirow{3}{*}{Gisette}
  & $10^{-1}$  & $.982 \pm .0003$ & 35.4
               & $.982 \pm .0003$ & 814.0
               & $\approx$ & $23.0\times$ \\
  & $10^{-3}$  & $.982 \pm .0002$ & 49.8
               & $.983 \pm .0001$ & 801.3
               & $\approx$ & $16.1\times$ \\
  & $10^{-5}$  & $.982 \pm .0002$ & 61.5
               & $.982 \pm .0002$ & 824.1
               & $\approx$ & $13.4\times$ \\
\midrule
\multirow{3}{*}{Dexter}
  & $10^{-1}$  & $.935 \pm .002$ & 21.0
               & $.934 \pm .002$ & 165.6
               & $\approx$ & $7.9\times$ \\
  & $10^{-3}$  & $.936 \pm .002$ & 29.2
               & $.936 \pm .002$ & 185.6
               & $\approx$ & $6.4\times$ \\
  & $10^{-5}$  & $.937 \pm .002$ & 34.1
               & $.936 \pm .002$ & 212.9
               & $\approx$ & $6.2\times$ \\
\midrule
\multirow{3}{*}{Riboflavin}
  & $10^{-1}$  & $.581 \pm .019$ & 47.6
               & $.539 \pm .020$ & 100.6
               & $\approx$ & $2.1\times$ \\
  & $10^{-3}$  & $.584 \pm .017$ & 64.6
               & $.542 \pm .021$ & 125.9
               & $\approx$ & $1.9\times$ \\
  & $10^{-5}$  & $.583 \pm .017$ & 86.8
               & $.545 \pm .021$ & 154.5
               & $\approx$ & $1.8\times$ \\
\midrule
\multirow{3}{*}{Isolet}
  & $10^{-1}$  & $.959 \pm .0003$ & 5.1
               & $.959 \pm .0002$ & 340.2
               & $\approx$ & $66.2\times$ \\
  & $10^{-3}$  & $.959 \pm .0003$ & 7.3
               & $.959 \pm .0003$ & 355.0
               & $\approx$ & $48.8\times$ \\
  & $10^{-5}$  & $.959 \pm .0003$ & 9.3
               & $.958 \pm .0002$ & 393.6
               & $\approx$ & $42.5\times$ \\
\bottomrule
\end{tabular}
\end{table}
 
Of the 18 (dataset, $\alpha$) configurations, 15 show no statistically significant difference in downstream performance between the two methods. In 3 configurations (Madelon at all three $\alpha$ levels), Vanilla Boruta achieves a statistically significant but practically small advantage: the largest observed gap is ${\sim}0.01$ in F1 (Madelon, $\alpha = 10^{-5}$).

Importantly, these performance differences, whether favourable or unfavourable to Greedy Boruta, are negligible in absolute terms, while the runtime reductions are substantial. Speedups range from $2.2\times$ (Riboflavin, $\alpha = 10^{-5}$) to $66.2\times$ (Isolet, $\alpha = 0.1$), with larger speedups observed on datasets where Vanilla Boruta requires more iterations to resolve features or trigger early-stopping, which depends on the signal-to-noise characteristics of the data rather than dimensionality alone. 

Figure~\ref{fig:realworld_boxplots} shows the full score distributions across seeds, confirming that the inter-method differences are small relative to the run-to-run variability within each method.

\begin{figure}[H]
    \centering
    \includegraphics[width=\textwidth, height=1\textheight, keepaspectratio]{score_comparison.png}
    \caption{Downstream performance (weighted F1 or $R^2$) across significance levels for each dataset (5 seeds per configuration). The narrow interquartile ranges and overlapping boxes confirm that performance differences between methods are small relative to run-to-run variability.}
    \label{fig:realworld_boxplots}
\end{figure}

\subsubsection{Pareto Efficiency}
\label{sec:experiments:realworld:pareto}

Figure~\ref{fig:pareto} plots mean downstream score against mean wall-clock time for all configurations. Across all six datasets, the Greedy Boruta points sit at the same vertical position as Vanilla Boruta but are displaced one to two orders of magnitude to the left on the time axis, Pareto-dominating or matching Vanilla Boruta in every configuration.

\begin{figure}
    \centering
    \includegraphics[width=\textwidth, height=0.42\textheight, keepaspectratio]{pareto_efficiency.png}
    \caption{Pareto efficiency: downstream score vs.\ wall-clock time (log scale). Each colour represents a dataset; circles are Greedy Boruta, crosses are Vanilla Boruta. Different points of the same colour and marker correspond to different $\alpha$ levels.}
    \label{fig:pareto}
\end{figure}

\subsubsection{Convergence and Early-Stopping Risk}
\label{sec:experiments:convergence}

To quantify the risk a practitioner faces when stopping Vanilla Boruta before convergence, we checkpoint Vanilla Boruta at every iteration on Gisette ($p = 5{,}000$, $\alpha = 0.1$) and track two metrics: the \emph{Undecidedness factor} $U = N_{\text{tentative}} / N_{\text{total}}$ (fraction of unresolved features) and the \emph{Potential miss rate} $M = N_{\text{tentative}} / (N_{\text{tentative}} + N_{\text{confirmed}})$ (fraction of possibly relevant features that would be lost if tentatives were discarded at that iteration).

\begin{figure}[H]
    \centering
    \includegraphics[width=\textwidth, height=0.28\textheight, keepaspectratio]{undecidedness_single_alpha_Gisette_0.1.png}
    \caption{Undecidedness factor (left) and Potential miss rate (right) for Vanilla Boruta on Gisette with $\alpha = 0.1$ as a function of iteration number. The dashed lines mark complete resolution (0\%).}
    \label{fig:undecidedness_gisette}
\end{figure}

Figure~\ref{fig:undecidedness_gisette} shows that after the full 512-iteration budget, 2--3\% of features remain tentative across all seeds, corresponding to 100--200 features resolved by the rough fix heuristic rather than the binomial test. The potential miss rate plateaus at approximately 7\%, meaning a practitioner who stops early risks discarding over 7\% of features that might eventually be confirmed. Greedy Boruta eliminates this uncertainty: all 5{,}000 features are resolved in $K = \lceil \log_2(5000/0.1) \rceil = 16$ iterations with no tentatives and no reliance on heuristic post-processing.

\section{Discussion}
\label{sec:discussion}

\subsection{The FPR--Convergence Trade-off in Practice}
\label{sec:discussion:tradeoff}

The theoretical analysis in Section~\ref{sec:method:fpr} establishes, under those assumptions, that Greedy Boruta exchanges Vanilla Boruta's near-zero FPR for a bounded FPR of $K/(p+1)$, where $K = \lceil \log_2(p/\alpha) \rceil$. The empirical results allow us to assess whether this trade-off is acceptable in practice. 

On the synthetic benchmark, Greedy Boruta incorrectly (and weakly) confirms between 4 and 11 noise features out of 395 across all seeds and $\alpha$ levels, all of them under the expected false positive rates under the estabilished FPR bounds (12.97 for $\alpha=10^{-1}$ and 25.94 for $\alpha=10^{-5}$, monotonically increasing).
% a small fraction of the noise space that produces no statistically significant degradation in F1 (Section~\ref{sec:experiments:alpha}). 
On the six real-world benchmarks, the only statistically significant differences occur on Madelon (all three $\alpha$ levels), with a maximum gap of ${\sim}0.01$ in F1 (Table~\ref{tab:realworld_results}). Madelon has the smallest feature count in our benchmark suite ($p = 500$), which is precisely the regime where the FPR bound from Theorem~\ref{thm:fpr} is least favorable. At larger $p$, the bound tightens: for $p = 5{,}000$ (Gisette) or $p = 20{,}000$ (Dexter), the per-feature expected FPR falls below $1\%$, and no significant performance difference is observed. This highlights the strength of Greedy Boruta on high-dimensional datasets, where the computational gains are more relevant and the FP rates are smallest. On the other hand, when dealing with low-dimensional datasets Vanilla Boruta may still be preferred since it may still be computationally reasonable, and provides lower FPR's. 

From a broader statistical perspective, Greedy Boruta can be understood as adopting a power-maximizing philosophy within the all-relevant framework. It sacrifices Type~I error control (from near-zero to bounded) in exchange for maximum recall, mirroring the well-known tension between family-wise error rate control and false discovery rate control in the multiple testing literature. The justification follows from the asymmetry discussed in Section~\ref{sec:introduction}: in the all-relevant paradigm, false negatives are irrecoverable while false positives can be filtered downstream.

\subsection{Limitations}
\label{sec:discussion:limitations}

\begin{enumerate}
    \item The FPR analysis relies on the shadow exchangeability assumption (Assumption~\ref{asm:exchangeability}), which requires that irrelevant features and shadow features be identically distributed with respect to the importance measure. Violations of this assumption, for instance due to importance measure biases toward high-cardinality features, affect both methods. However, features with marginal biases that produce occasional but non-significant hits are weakly confirmed by Greedy Boruta but may remain tentative or be rejected under Vanilla Boruta's stricter criterion.

    \item The parallelization benefit depends on the base estimator's internal parallelism; when tree training saturates all CPU cores, the dominant speedup comes from reduced iteration counts rather than outer parallelism (Section~\ref{sec:method:parallel}).

    \item Finally, Greedy Boruta inherits all base estimator limitations from Vanilla Boruta, including sensitivity to the choice of importance measure and the assumption that tree-based ensembles adequately capture feature relevance.
\end{enumerate}

% \section{Conclusion}
% \label{sec:conclusion}

% We have presented Greedy Boruta, a modification to the Boruta algorithm that replaces the binomial test confirmation criterion with a single-hit rule: a feature is confirmed if it ever exceeds the maximum shadow importance in any iteration. This change yields three formal guarantees: a true positive rate (feature recall) at least as high as Vanilla Boruta (Theorem~\ref{thm:recall}), convergence in exactly $K = \lceil \log_2(p/\alpha) \rceil$ iterations under Bonferroni correction (Theorem~\ref{thm:convergence}) and a per-feature false positive rate bounded above by $K/(p+1)$ (Theorem~\ref{thm:fpr}). A further consequence is that FDR correction becomes trivially redundant (Theorem~\ref{thm:fdr_trivial}), collapsing the multiple testing correction landscape to a single Boolean parameter. Experiments on a synthetic benchmark with known ground truth and six real-world high-dimensional datasets confirm these theoretical properties and demonstrate substantial speedups with negligible impact on downstream model performance.

% We recommend Greedy Boruta as the default choice for all-relevant feature selection in pipelines where downstream filtering is available, particularly in high-dimensional settings where its FPR bound is tightest and its computational advantage is largest. An implementation is publicly available as \texttt{greedyboruta} in PyPI. Future work includes extending the greedy criterion to alternative importance measures such as SHAP and permutation importance and investigating adaptive significance level selection strategies.

\section{Conclusion}
\label{sec:conclusion}

We have presented Greedy Boruta, a modification to the Boruta algorithm that replaces the binomial-test confirmation criterion with a single-hit rule: a feature is confirmed if it exceeds the maximum shadow importance in at least one iteration. This modification yields four formal guarantees: a feature recall at least as high as Vanilla Boruta (Theorem~\ref{thm:recall}), convergence in exactly $K=\lceil \log_2(p/\alpha)\rceil$ iterations under Bonferroni correction (Theorem~\ref{thm:convergence}), a per-feature false positive rate bounded above by $K/(p+1)$ (Theorem~\ref{thm:fpr}) and the triviality of FDR correction under the greedy criterion (Theorem~\ref{thm:fdr_trivial}). As a consequence, Greedy Boruta eliminates all termination-related hyperparameters and substantially simplifies the multiple-testing correction framework used by Boruta.

Beyond the core algorithm, we introduced two practical extensions. First, we proposed a weak/strong confirmation framework that distinguishes features satisfying the original Boruta confirmation criterion from those selected exclusively by the greedy criterion, providing users with an interpretable measure of confidence among the confirmed features. Second, in Appendix \ref{app:add_results_reducing_shadows}, we showed that the number of shadow features can be reduced from $p$ to $r \leq p$, yielding an explicit trade-off between computational efficiency and false positive control. This extension is particularly attractive in high-dimensional and distributed computing environments, where the creation and storage of shadow features can become a substantial computational bottleneck.

Experiments on a synthetic benchmark with known ground truth and six real-world high-dimensional datasets support the theoretical analysis and demonstrate substantial computational gains, with speedups ranging from $1.8$ to $66.2\times$ while maintaining downstream predictive performance that is statistically indistinguishable from Vanilla Boruta in most configurations.

We recommend Greedy Boruta as a practical default for all-relevant feature selection when downstream filtering, regularization or model validation procedures are available. The method is particularly appealing in high-dimensional settings, where its false positive rate remains small while its computational advantages become most pronounced. Future work includes extending the greedy criterion to alternative feature-importance measures, such as SHAP and permutation importance, and investigating the statistical properties of weakly confirmed features in large-scale applications. An open-source implementation is publicly available through the \texttt{greedyboruta} package on PyPI.

% =============================================================================
% ACKNOWLEDGMENTS (for camera-ready only)
% =============================================================================
% \paragraph{Acknowledgments.} We thank [names] for helpful discussions.

% =============================================================================
% REFERENCES
% =============================================================================

\bibliographystyle{splncs04}
\bibliography{references}

\appendix
\renewcommand{\theHsection}{\Alph{section}}

\section{Proofs}
\label{app:proofs}

For each feature $j$, define the total number of hits across $K$ iterations as
\[
H_j = \sum_{i=1}^{K} \mathbf{1}\!\left\{ I_j^{(i)} > M^{(i)} \right\},
\]
where $I_j^{(i)}$ denotes the importance of feature $j$ at iteration $i$ and $M^{(i)} = \max_j S_j^{(i)}$ denotes the maximum shadow importance at iteration $i$.

\subsection{Proof of Theorem~\ref{thm:recall} (Recall Guarantee)}

\begin{proof}
Vanilla Boruta confirms features through two mechanisms: the sequential binomial test during execution and, upon termination, the rough fix heuristic applied to remaining tentative features. We show that both mechanisms require $H_j \geq 1$.


\paragraph{Case 1: Binomial test confirmation.}
Consider any feature $j$ confirmed by the binomial test, which tests 
$$
H_0 : p_j \leq 0.5
\quad \text{versus} \quad
H_1 : p_j > 0.5.
$$
This requires $\mathbb{P}(X \geq H_j \mid n = K, p = 0.5) < \alpha$, where $X \sim \text{Binomial}(K, 0.5)$ and $\alpha$ denotes the significance level. For $H_j = 0$, this probability equals $\mathbb{P}(X \geq 0) = 1 > \alpha$ for any $\alpha < 1$, so the hypothesis test cannot reject $H_0$ and a feature with $H_j = 0$ is never confirmed. Therefore, for a feature to be confirmed as relevant, $H_j \geq 1$.

\paragraph{Case 2: Rough fix confirmation.}
The rough fix~\cite{kursa2010feature} confirms a tentative feature $j$ if its median importance exceeds the median of the per-iteration maximum shadow importances, i.e., 
$$\mathrm{median}(I_j^{(1)}, \ldots, I_j^{(K)}) > \mathrm{median}(M^{(1)}, \ldots, M^{(K)}).$$

% where $f_i$ is the importance of feature $f$ in iteration $i$ and $M_i = \max_j S_j^{(i)}$ is the maximum shadow importance.

If $H_j = 0$, then $I_j^{(i)} \leq M^{(i)}$ for every iteration $i$. The empirical distribution of the importance of feature $j$ is therefore stochastically dominated (see \cite{levy1990stochastic,blackwell1953equivalent}) by that of the maximum shadow importances. Stochastic dominance preserves quantile ordering, so $\mathrm{median}(I_j^{(1)}, \ldots, I_j^{(K)}) \leq \mathrm{median}(M^{(1)}, \ldots, M^{(K)})$, and the rough fix cannot confirm a feature with zero hits.

In both cases, confirmation by Vanilla Boruta requires $H_j \geq 1$ for all $j$. By the Greedy Boruta confirmation criterion (Algorithm~\ref{alg:greedy_boruta}, line 14), any such feature is also confirmed, so $j \in \mathcal{C}_G$. Since this holds for all $j \in \mathcal{C}_V$, we conclude $\mathcal{C}_V \subseteq \mathcal{C}_G$. 
\end{proof}

\subsection{Proof of Theorem~\ref{thm:convergence} (Convergence Guarantee)}

\begin{proof}
At iteration $K$, partition the features into
$$
\mathcal{H}^{+} = \{j : H_j \geq 1\}
\quad\text{and}\quad
\mathcal{H}^{0} = \{j : H_j = 0\}.
$$
By Algorithm~\ref{alg:greedy_boruta} (line~14), all features in $\mathcal{H}^{+}$ are confirmed by Greedy Boruta. In contrast, for any feature $j \in \mathcal{H}^{0}$, the rejection test evaluates
$$
H_0 : p_j \geq 0.5
\quad \text{versus} \quad
H_1 : p_j < 0.5,
$$
where $p_j$ denotes the probability that feature $j$ exceeds the maximum shadow importance in a given iteration. Under the null hypothesis, the hit count follows
$$
H_j \sim \mathrm{Binomial}(K,0.5).
$$
Since $H_j = 0$, the rejection p-value is
$$
\mathbb{P}(H_j = 0 \mid K,0.5)
=
(0.5)^K
=
2^{-K}.
$$
Notice that all features in $\mathcal{H}^0$ have this identical p-value. Under Bonferroni correction across $|\mathcal{H}^0| \leq p$ features, rejection requires $2^{-K} < \alpha/p$. The minimum integer satisfying this is $K = \lceil \log_2(p/\alpha) \rceil$. In practice, at iteration $K$ all features in $\mathcal{H}^0$ are rejected and all in $\mathcal{H}^+$ are confirmed, so the algorithm terminates with no tentative features.
\end{proof}
\noindent
\textit{Note}: Without Bonferroni correction, rejection requires $2^{-K}  < \alpha$, giving $K = \lceil \log_2(1/\alpha) \rceil$. 

\subsection{Proof of Theorem~\ref{thm:fdr_trivial} (FDR Correction Triviality)}

\begin{proof}
From the proof of Theorem~\ref{thm:convergence}, any tentative feature $j$ at iteration $K$ has $H_j = 0$ and p-value $p^* = 2^{-K}$.

The Benjamini-Hochberg procedure orders p-values $p_{(1)} \leq \cdots \leq p_{(m)}$ and finds the largest $i$ such that $p_{(i)} \leq (i/m) \cdot \alpha$. Since all p-values are identical ($p_{(i)} = p^*$ for all $i$), this reduces to checking whether $p^* \leq \alpha$. From the construction of $K = \lceil \log_2(1/\alpha) \rceil$, we have $p^* = 2^{-K} < \alpha$, so the condition holds at $i = m$ and all tentative features are rejected. Hence, the FDR-corrected and uncorrected procedures coincide.
\end{proof}

\subsection{Proof of Theorem~\ref{thm:fpr} (FPR of Greedy Boruta)}

\begin{proof}
Under Assumption~\ref{asm:exchangeability}, the importance of an irrelevant feature and the $p$ shadow importances form $p+1$ exchangeable random variables. By symmetry, each variable is equally likely to be the maximum importance, so the per-iteration probability that an irrelevant feature receives a hit is
$$
\theta_0 = \frac{1}{p+1}.
$$
Since the iterations are independent, $H_j \sim \mathrm{Binomial}(K,\theta_0)$, and under the greedy criterion, a false positive occurs whenever an irrelevant feature obtains at least one hit. Therefore, the false positive probability is
\begin{align}
\mathbb{P}(H_j \geq 1)
&= 1 - \mathbb{P}(H_j = 0) \nonumber \\
&= 1 - (1-\theta_0)^K \nonumber \\
&= 1 - \left(\frac{p}{p+1}\right)^K. \nonumber
\end{align}
For the upper bound, Bernoulli's inequality~\cite{mitrinovic1970analytic} states that $(1-x)^K \geq 1-Kx$, for $x \in [0,1]$ and integer $K \geq 1$. Setting $x = 1/(p+1)$ yields
$$
\left(\frac{p}{p+1}\right)^K \geq 1 - \frac{K}{p+1}.
$$
Substituting into the exact expression above gives
$$
\mathbb{P}(H_j \geq 1)
\leq \frac{K}{p+1}.
$$
\end{proof}

\paragraph{Comparison with Vanilla Boruta.}
Under Assumption~\ref{asm:exchangeability}, the FPR of Vanilla Boruta is effectively negligible. In Vanilla Boruta, a feature is confirmed only if its hit count exceeds a threshold \(h^*\) satisfying
\[
\mathbb{P}(\mathrm{Binomial}(K,0.5) \geq h^*) \leq \alpha.
\]
For an irrelevant feature, the per-iteration hit probability is \(\theta_0 = 1/(p+1)\), so its hit count follows
\[
H_j \sim \mathrm{Binomial}(K,\theta_0).
\]
Consequently, the probability that an irrelevant feature is falsely confirmed by Vanilla Boruta is $\mathbb{P}(H_j \geq h^*)$, which decays exponentially in \(K\) and is therefore extremely small in high-dimensional settings. For example, with \(\alpha = 0.001\) and \(p = 500\), Greedy Boruta requires \(K = 19\) iterations, while Vanilla Boruta requires at least \(h^* = 17\) hits for confirmation. Under the null hit probability \(\theta_0 = 1/501\), the resulting false positive probability for Vanilla Boruta is approximately \(10^{-44}\), compared to the false positive probability under Greedy Boruta
\[
\mathbb{P}(H_j \geq 1)
=
1 - \left(\frac{500}{501}\right)^{19}
\approx 3.73\%.
\]
This contrast reflects the central trade-off between the two procedures: Greedy Boruta sacrifices the near-zero FPR of Vanilla Boruta in exchange for deterministic convergence, elimination of tentative features, and substantially improved computational efficiency.

% \subsection{Proof of Proposition~\ref{prop:mr_mp}}

% \begin{proof}
% The event \(\{M_r = M_p\}\) occurs if, and only if, all of the remaining \(p-r\) variables are smaller than \(M_r\). Conditioning on \(M_r\),
% \[
% \mathbb{P}(M_r = M_p \mid M_r)
% =
% F(M_r)^{p-r}.
% \]
% Taking the expectation yields
% \[
% \mathbb{P}(M_r = M_p)
% =
% \mathbb{E}\!\left[F(M_r)^{p-r}\right].
% \]
% We now derive the distribution of \(F(M_r)\). For any \(u \in [0,1]\),
% \begin{align*}
% \mathbb{P}(F(M_r) \leq u)
% &=
% \mathbb{P}(M_r \leq F^{-1}(u)) \\
% &=
% \mathbb{P}(S_1 \leq F^{-1}(u),\ldots,S_r \leq F^{-1}(u)) \\
% &=
% \prod_{i=1}^{r} \mathbb{P}(S_i \leq F^{-1}(u)) \\
% &=
% u^r.
% \end{align*}
% Therefore,
% \[
% F(M_r) \sim \mathrm{Beta}(r,1),
% \]
% with density
% \[
% f(u)=r\,u^{r-1},
% \qquad 0 \leq u \leq 1.
% \]
% Consequently,
% \begin{align*}
% \mathbb{P}(M_r = M_p)
% &=
% \mathbb{E}\!\left[F(M_r)^{p-r}\right] \\
% &=
% \int_0^1 u^{p-r} \, r u^{r-1}\,du \\
% &=
% r\int_0^1 u^{p-1}\,du \\
% &=
% \frac{r}{p}.
% \end{align*}
% \end{proof}

\section{Additional results}
\label{app:add_results_reducing_shadows}

\subsection{FPR analysis with $r \leq p$ shadow features}
\begin{figure}[H]
    \centering
    \includegraphics[width=\linewidth]{FPR_with_shadows.png}
    \caption{Upper bound on per-feature false positive rate (FPR) of Greedy Boruta under Bonferroni correction using different number of features, fraction of shadow features and significance level.}
    \label{fig:fpr_shadows}
\end{figure}

\begin{table}[H]
\centering
\caption{Upper bound on per-feature false positive rate (FPR) of Greedy Boruta under Bonferroni correction using different number of features, fraction of shadow features and significance level. The upper bound on FPR is computed as $K/(r+1)$, where $K=\lceil \log_2(p/\alpha)\rceil$.}
\begin{subtable}{\textwidth}
\centering
\caption{$r = 0.10p$}
\begin{tabular}{c|cc|cc|cc|cc}
\toprule
& \multicolumn{2}{c|}{$p = 100$}
& \multicolumn{2}{c|}{$p = 500$}
& \multicolumn{2}{c|}{$p = 1{,}000$}
& \multicolumn{2}{c}{$p = 10{,}000$} \\
$\alpha$
& $K$ & FPR
& $K$ & FPR
& $K$ & FPR
& $K$ & FPR \\
\midrule
$0.05$    & 11 & 100.00\% & 14 & 27.45\% & 15 & 14.85\% & 18 & 1.80\% \\
$0.01$    & 14 & 100.00\% & 16 & 31.37\% & 17 & 16.83\% & 20 & 2.00\% \\
$0.001$   & 17 & 100.00\% & 19 & 37.25\% & 20 & 19.80\% & 24 & 2.40\% \\
$0.0001$  & 20 & 100.00\% & 23 & 45.10\% & 24 & 23.76\% & 27 & 2.70\% \\
$0.00001$ & 24 & 100.00\% & 26 & 50.98\% & 27 & 26.73\% & 30 & 3.00\% \\
\bottomrule
\end{tabular}
\end{subtable}

% \vspace{0.1cm}

\begin{subtable}{\textwidth}
\centering
\caption{$r = 0.25p$}
\begin{tabular}{c|cc|cc|cc|cc}
\toprule
& \multicolumn{2}{c|}{$p = 100$}
& \multicolumn{2}{c|}{$p = 500$}
& \multicolumn{2}{c|}{$p = 1{,}000$}
& \multicolumn{2}{c}{$p = 10{,}000$} \\
$\alpha$
& $K$ & FPR
& $K$ & FPR
& $K$ & FPR
& $K$ & FPR \\
\midrule
$0.05$    & 11 & 42.31\% & 14 & 11.02\% & 15 & 5.98\% & 18 & 0.72\% \\
$0.01$    & 14 & 53.85\% & 16 & 12.60\% & 17 & 6.77\% & 20 & 0.80\% \\
$0.001$   & 17 & 65.38\% & 19 & 14.96\% & 20 & 7.97\% & 24 & 0.96\% \\
$0.0001$  & 20 & 76.92\% & 23 & 18.11\% & 24 & 9.56\% & 27 & 1.08\% \\
$0.00001$ & 24 & 92.31\% & 26 & 20.47\% & 27 & 10.76\% & 30 & 1.20\% \\
\bottomrule
\end{tabular}
\end{subtable}

% \vspace{0.1cm}

\begin{subtable}{\textwidth}
\centering
\caption{$r = 0.50p$}
\begin{tabular}{c|cc|cc|cc|cc}
\toprule
& \multicolumn{2}{c|}{$p = 100$}
& \multicolumn{2}{c|}{$p = 500$}
& \multicolumn{2}{c|}{$p = 1{,}000$}
& \multicolumn{2}{c}{$p = 10{,}000$} \\
$\alpha$
& $K$ & FPR
& $K$ & FPR
& $K$ & FPR
& $K$ & FPR \\
\midrule
$0.05$    & 11 & 21.57\% & 14 & 5.58\% & 15 & 2.99\% & 18 & 0.36\% \\
$0.01$    & 14 & 27.45\% & 16 & 6.37\% & 17 & 3.39\% & 20 & 0.40\% \\
$0.001$   & 17 & 33.33\% & 19 & 7.57\% & 20 & 3.99\% & 24 & 0.48\% \\
$0.0001$  & 20 & 39.22\% & 23 & 9.16\% & 24 & 4.79\% & 27 & 0.54\% \\
$0.00001$ & 24 & 47.06\% & 26 & 10.36\% & 27 & 5.39\% & 30 & 0.60\% \\
\bottomrule
\end{tabular}
\end{subtable}

% \vspace{0.1cm}

\begin{subtable}{\textwidth}
\centering
\caption{$r = 0.75p$}
\begin{tabular}{c|cc|cc|cc|cc}
\toprule
& \multicolumn{2}{c|}{$p = 100$}
& \multicolumn{2}{c|}{$p = 500$}
& \multicolumn{2}{c|}{$p = 1{,}000$}
& \multicolumn{2}{c}{$p = 10{,}000$} \\
$\alpha$
& $K$ & FPR
& $K$ & FPR
& $K$ & FPR
& $K$ & FPR \\
\midrule
$0.05$    & 11 & 14.47\% & 14 & 3.72\% & 15 & 2.00\% & 18 & 0.24\% \\
$0.01$    & 14 & 18.42\% & 16 & 4.25\% & 17 & 2.26\% & 20 & 0.27\% \\
$0.001$   & 17 & 22.37\% & 19 & 5.05\% & 20 & 2.66\% & 24 & 0.32\% \\
$0.0001$  & 20 & 26.32\% & 23 & 6.12\% & 24 & 3.19\% & 27 & 0.36\% \\
$0.00001$ & 24 & 31.58\% & 26 & 6.91\% & 27 & 3.59\% & 30 & 0.40\% \\
\bottomrule
\end{tabular}
\end{subtable}

% \vspace{0.1cm}

\begin{subtable}{\textwidth}
\centering
\caption{$r = p$ (original Greedy Boruta)}
\begin{tabular}{c|cc|cc|cc|cc}
\toprule
& \multicolumn{2}{c|}{$p = 100$}
& \multicolumn{2}{c|}{$p = 500$}
& \multicolumn{2}{c|}{$p = 1{,}000$}
& \multicolumn{2}{c}{$p = 10{,}000$} \\
$\alpha$
& $K$ & FPR
& $K$ & FPR
& $K$ & FPR
& $K$ & FPR \\
\midrule
$0.05$    & 11 & 10.89\% & 14 & 2.79\% & 15 & 1.50\% & 18 & 0.18\% \\
$0.01$    & 14 & 13.86\% & 16 & 3.19\% & 17 & 1.70\% & 20 & 0.20\% \\
$0.001$   & 17 & 16.83\% & 19 & 3.79\% & 20 & 2.00\% & 24 & 0.24\% \\
$0.0001$  & 20 & 19.80\% & 23 & 4.59\% & 24 & 2.40\% & 27 & 0.27\% \\
$0.00001$ & 24 & 23.76\% & 26 & 5.19\% & 27 & 2.70\% & 30 & 0.30\% \\
\bottomrule
\end{tabular}
\end{subtable}
\label{tab:fpr_shadows}
\end{table}

\begin{table}[H]
\centering
\caption{Exact per-feature false positive rate (FPR) of Greedy Boruta under Bonferroni correction using different numbers of features, fractions of shadow features, and significance levels. The FPR is computed as $1-(1-\frac{1}{r+1})^K$, where $K=\lceil \log_2(p/\alpha)\rceil$.}
\label{tab:fpr_shadows_exact}

\begin{subtable}{\textwidth}
\centering
\caption{$r = 0.10p$}
\begin{tabular}{c|cc|cc|cc|cc}
\toprule
& \multicolumn{2}{c|}{$p = 100$}
& \multicolumn{2}{c|}{$p = 500$}
& \multicolumn{2}{c|}{$p = 1{,}000$}
& \multicolumn{2}{c}{$p = 10{,}000$} \\
$\alpha$
& $K$ & FPR
& $K$ & FPR
& $K$ & FPR
& $K$ & FPR \\
\midrule
$0.05$    & 11 & 64.95\% & 14 & 24.21\% & 15 & 13.87\% & 18 & 1.78\% \\
$0.01$    & 14 & 73.67\% & 16 & 27.16\% & 17 & 15.56\% & 20 & 1.98\% \\
$0.001$   & 17 & 80.22\% & 19 & 31.36\% & 20 & 18.05\% & 24 & 2.37\% \\
$0.0001$  & 20 & 85.14\% & 23 & 36.58\% & 24 & 21.24\% & 27 & 2.66\% \\
$0.00001$ & 24 & 89.85\% & 26 & 40.24\% & 27 & 23.56\% & 30 & 2.95\% \\
\bottomrule
\end{tabular}
\end{subtable}

% \vspace{0.3cm}

\begin{subtable}{\textwidth}
\centering
\caption{$r = 0.25p$}
\begin{tabular}{c|cc|cc|cc|cc}
\toprule
& \multicolumn{2}{c|}{$p = 100$}
& \multicolumn{2}{c|}{$p = 500$}
& \multicolumn{2}{c|}{$p = 1{,}000$}
& \multicolumn{2}{c}{$p = 10{,}000$} \\
$\alpha$
& $K$ & FPR
& $K$ & FPR
& $K$ & FPR
& $K$ & FPR \\
\midrule
$0.05$    & 11 & 35.04\% & 14 & 10.56\% & 15 & 5.81\% & 18 & 0.72\% \\
$0.01$    & 14 & 42.25\% & 16 & 11.97\% & 17 & 6.56\% & 20 & 0.80\% \\
$0.001$   & 17 & 48.66\% & 19 & 14.05\% & 20 & 7.67\% & 24 & 0.96\% \\
$0.0001$  & 20 & 54.36\% & 23 & 16.75\% & 24 & 9.14\% & 27 & 1.07\% \\
$0.00001$ & 24 & 60.99\% & 26 & 18.71\% & 27 & 10.22\% & 30 & 1.19\% \\
\bottomrule
\end{tabular}
\end{subtable}

% \vspace{0.3cm}

\begin{subtable}{\textwidth}
\centering
\caption{$r = 0.50p$}
\begin{tabular}{c|cc|cc|cc|cc}
\toprule
& \multicolumn{2}{c|}{$p = 100$}
& \multicolumn{2}{c|}{$p = 500$}
& \multicolumn{2}{c|}{$p = 1{,}000$}
& \multicolumn{2}{c}{$p = 10{,}000$} \\
$\alpha$
& $K$ & FPR
& $K$ & FPR
& $K$ & FPR
& $K$ & FPR \\
\midrule
$0.05$    & 11 & 19.57\% & 14 & 5.44\% & 15 & 2.95\% & 18 & 0.36\% \\
$0.01$    & 14 & 24.21\% & 16 & 6.19\% & 17 & 3.34\% & 20 & 0.40\% \\
$0.001$   & 17 & 28.58\% & 19 & 7.30\% & 20 & 3.92\% & 24 & 0.48\% \\
$0.0001$  & 20 & 32.70\% & 23 & 8.77\% & 24 & 4.68\% & 27 & 0.54\% \\
$0.00001$ & 24 & 37.83\% & 26 & 9.86\% & 27 & 5.25\% & 30 & 0.60\% \\
\bottomrule
\end{tabular}
\end{subtable}

% \vspace{0.3cm}

\begin{subtable}{\textwidth}
\centering
\caption{$r = 0.75p$}
\begin{tabular}{c|cc|cc|cc|cc}
\toprule
& \multicolumn{2}{c|}{$p = 100$}
& \multicolumn{2}{c|}{$p = 500$}
& \multicolumn{2}{c|}{$p = 1{,}000$}
& \multicolumn{2}{c}{$p = 10{,}000$} \\
$\alpha$
& $K$ & FPR
& $K$ & FPR
& $K$ & FPR
& $K$ & FPR \\
\midrule
$0.05$    & 11 & 13.56\% & 14 & 3.66\% & 15 & 1.98\% & 18 & 0.24\% \\
$0.01$    & 14 & 16.93\% & 16 & 4.17\% & 17 & 2.24\% & 20 & 0.27\% \\
$0.001$   & 17 & 20.16\% & 19 & 4.93\% & 20 & 2.63\% & 24 & 0.32\% \\
$0.0001$  & 20 & 23.27\% & 23 & 5.94\% & 24 & 3.15\% & 27 & 0.36\% \\
$0.00001$ & 24 & 27.23\% & 26 & 6.69\% & 27 & 3.53\% & 30 & 0.40\% \\
\bottomrule
\end{tabular}
\end{subtable}

% \vspace{0.3cm}

\begin{subtable}{\textwidth}
\centering
\caption{$r = p$ (original Greedy Boruta)}
\begin{tabular}{c|cc|cc|cc|cc}
\toprule
& \multicolumn{2}{c|}{$p = 100$}
& \multicolumn{2}{c|}{$p = 500$}
& \multicolumn{2}{c|}{$p = 1{,}000$}
& \multicolumn{2}{c}{$p = 10{,}000$} \\
$\alpha$
& $K$ & FPR
& $K$ & FPR
& $K$ & FPR
& $K$ & FPR \\
\midrule
$0.05$    & 11 & 10.37\% & 14 & 2.76\% & 15 & 1.49\% & 18 & 0.18\% \\
$0.01$    & 14 & 13.00\% & 16 & 3.15\% & 17 & 1.68\% & 20 & 0.20\% \\
$0.001$   & 17 & 15.56\% & 19 & 3.73\% & 20 & 1.98\% & 24 & 0.24\% \\
$0.0001$  & 20 & 18.05\% & 23 & 4.49\% & 24 & 2.37\% & 27 & 0.27\% \\
$0.00001$ & 24 & 21.24\% & 26 & 5.06\% & 27 & 2.66\% & 30 & 0.30\% \\
\bottomrule
\end{tabular}
\end{subtable}

\end{table}

\subsection{Reduced-Shadow Greedy Boruta on Real-World Benchmarks}

\end{document}

